\section{Discussion and Conclusions}

	Several approaches for spatiotemporal shape modeling of anatomical structures were discussed in this review. Rapid developments in the field, especially in recent years, have been fueled by advancements in DL and are set to only continually progress further. Nevertheless, the works found in the existing literature have been mainly focused on incremental developments in methodology or applications of novel new tools. This is in contrast with applying already developed tools to existing or novel clinical challenges. This seeming reluctance of the medical imaging community towards application-based research could stem from a multitude of reasons, but a simple lack of data could be the main factor, as we will discuss shortly. Deficiencies notwithstanding, in this section, we will discuss key concepts of spatiotemporal shape modeling uncovered from our review. We will then outline several key barriers to further research in the field before speculating on future research directions. 
	
	\subsection{Nonlinear shape manifolds} 

		From our review, it is clear that anatomical shape variation is highly nonlinear. This nonlinearity is further compounded by the additional nonlinear dynamics of growth and changing biological structures over time, leading to an intricate and complex outlook. Thus, the best-suited models for spatiotemporal progression are those that lie on non-Euclidean manifolds as they best capture this inherently high dimensional problem. A potential reason for this could be the manifold hypothesis, wherein it is postulated that all high-dimensional data lie on an embedded low-dimensional manifold \cite{Fefferman2016, NIPS2010_8a1e808b}. The task of spatiotemporal shape modeling can then be reduced to identifying and characterizing these manifolds, either implicitly or explicitly. The LDDMM framework discussed in Section \ref{sec:LDDMM}, for example, explicitly seeks to uncover spatiotemporal trajectories of diffeomorphisms traversing across a manifold. DL techniques discussed in Section \ref{sec:DL} also implicitly benefit from manifolds, as the efficacy of DL techniques has been attributed to their capability to uncover and disentangle underlying the manifolds of complex data \cite{Brahma2016}.    

		In contrast to manifold-based techniques, several works do exist that have attempted to extend linear (PCA-based) statistical shape models towards spatiotemporal shape models. These, however, fall short when compared to LDDMM and DL-based solutions as they effectively only serve to compare differences across and interpolate between time points as opposed to true longitudinal forecasting \cite{Hamarneh2004, Kasahara2018, BinteAlam2020, Saito2019}. Due to their reliance on landmarks, these methods do not effectively work if anatomies significantly change over time, as is the case, especially in early development. Furthermore, they are incapable of separating groupwise versus individual developmental trends, nor are they capable of effective data imputation \cite{Adams2023}. Therefore, while these methods might be effective for comparing shape variation across time points, they are not as effective for shape trajectory forecasting as manifold-based methods. 
	
		Comparatively, manifold-based techniques are more effective as the longitudinal trajectories traversing the shape space yield an effective description of shape variation over time. LDDMM techniques offer a structured framework to describe shape variation, and furthermore, the geodesic trajectories themselves are clinically relevant as they offer prognostic and diagnostic utility. When utilized within a hierarchical model that incorporates many trajectories for a population average, new trajectories can be estimated for unseen data, which could offer prognostic significance. Furthermore, trajectories can be compared using relational transport operators to diagnose if a trajectory is irregular compared to population averages. Similarly, DL methods mostly operate directly on medical imaging data with convolutional networks. This allows us a way to extract hidden features from images which could also influence spatiotemporal trajectories, otherwise lost during parameterization processes necessary for LDDMM or PCA-based models. In encoding networks especially, the latent space encompassed by these extracted latent variables can be structured to construct a physically meaningful underlying spatiotemporal manifold. The inductive capacity of DL methods with such structured latent spaces is then superior to linear methods, capable of imputing missing data and predicting spatiotemporal trajectories. 
			
	\subsection{Paucity of longitudinal datasets}	
			
		Another clear deficiency is the lack of large, open-source, and high-quality longitudinal imaging datasets. Existing datasets used in studies are generally small, in-house, cover a short time span, and are limited to very specific clinical conditions (Table \ref{tab:my-table}). This is, of course, understandable as it is extremely difficult to gather longitudinal data. Issues such as participant attrition \cite{Young2006} and ethical concerns \cite{Tinker2009} are just two examples of difficulties that hamper the execution of effective studies. An exception to this is the Alzheimer's Disease Neuroimaging Initiative (ADNI) database, which is a large multimodal database of longitudinal biomarker and neuroimaging data tracking the progression of AD \cite{Jack2008}. This dataset is particularly outstanding due to its size and comprehensiveness, leading to many studies covered in this review validating their methods on the ADNI dataset. Nevertheless, this dataset remains unique and standout compared to others. This paucity of longitudinal datasets, especially for medical imaging, impairs the efficacy of both LDDMM and DL techniques covered in this study. 
		
\begin{table}[]
\resizebox{\textwidth}{!}{
\begin{tabular}{lllccll}
\hline
\rowcolor[HTML]{EFEFEF} 
\multicolumn{1}{|c|}{\cellcolor[HTML]{EFEFEF}} &
  \multicolumn{1}{c|}{\cellcolor[HTML]{EFEFEF}} &
  \multicolumn{1}{c|}{\cellcolor[HTML]{EFEFEF}} &
  \multicolumn{2}{c|}{\cellcolor[HTML]{EFEFEF}\textbf{Age range}} & 
  \multicolumn{1}{c|}{\cellcolor[HTML]{EFEFEF}} &
  \multicolumn{1}{c|}{\cellcolor[HTML]{EFEFEF}} \\ \cline{4-5}
\rowcolor[HTML]{EFEFEF} 
\multicolumn{1}{|c|}{\multirow{-2}{*}{\cellcolor[HTML]{EFEFEF}\textbf{Name}}} &
  \multicolumn{1}{c|}{\multirow{-2}{*}{\cellcolor[HTML]{EFEFEF}\textbf{Anatomy}}} &
  \multicolumn{1}{c|}{\multirow{-2}{*}{\cellcolor[HTML]{EFEFEF}\textbf{Modality}}} &
  \multicolumn{1}{l|}{\cellcolor[HTML]{EFEFEF}\textbf{Youngest}} &
  \multicolumn{1}{l|}{\cellcolor[HTML]{EFEFEF}\textbf{Oldest}} &
  \multicolumn{1}{c|}{\multirow{-2}{*}{\cellcolor[HTML]{EFEFEF}\textbf{Number of subjects}}} &
  \multicolumn{1}{c|}{\multirow{-2}{*}{\cellcolor[HTML]{EFEFEF}\textbf{Repeated measurements}}} \\ \hline
OASIS-2                      & Brain      & MRI & 60                   & 96                   & 150                  & \textless{}=5                   \\ \hline
OASIS-3                      & Brain      & MRI & 42                   & 95                   & 1378                 & \textless{}=7                   \\ \hline
HABS-HD                      & Brain      & MRI & 40                   & 92                   & 3838                 & \textless{}=3                   \\ \hline
Harvard Aging Brain Study    & Brain      & MRI & 62                   & 90                   & \textgreater{}290    & \textless{}=3                   \\ \hline
ADNI-4                       & Brain      & MRI & 55                   & 90                   & \textgreater{}2400   & \textless{}=6                   \\ \hline
                             & Brain      & MRI &                      &                      & \textgreater{}5000   &                                 \\
                             & Whole body & DXA &                      &                      & \textgreater{}=5156  &                                 \\
                             & Abdomen    & MRI &                      &                      & \textgreater{}=11365 &                                 \\ 
\multirow{-4}{*}{UK Biobank} & Heart      & MRI & \multirow{-4}{*}{44} & \multirow{-4}{*}{87} & \textgreater{}=5100  & \multirow{-4}{*}{\textless{}=2} \\ \hline
MCSA                         & Brain      & MRI & 30                   & 89                   & 1802                 & ??? \\ \hline                           
\end{tabular}}
\vspace{0.2cm}
\caption{A summary of several longitudinal medical imaging datasets. OASIS - Open Access Series of Imaging Studies. HABS-HD - Health and Aging Brain Study - Health Disparities. ADNI - Alzheimer's Disease Neuroimaging Initiative. MCSA - Mayo Clinic Study of Aging.}
\label{tab:my-table}
\end{table}

		DL techniques are notoriously data hungry, with larger dataset sizes contributing significantly towards improved efficacy of networks \cite{Sun_2017_ICCV, JungKye15}. While techniques such as transfer learning \cite{Alzubaidi2020} and data augmentation \cite{Mumuni2022} seek to ameliorate this issue, it remains pervasive. Conversely, whilst the LDDMM framework is comparatively not as data-hungry, sufficiently sized datasets are also essential. Adequately sized and diverse datasets are vital to ensure that the estimated population average trajectories are reflective of the entire population. Solutions such as GAN-based frameworks discussed in Section 3.2 are shown to be helpful in addressing the issue of data paucity. Therein, generative processes and adversarial training frameworks can increase the generalizability of networks. The latter is particularly useful as the adversarial process assists in regularizing and structuring the latent space, implicitly learning the underlying spatiotemporal manifold. Nonetheless, the lack of datasets presents another issue of validity. In essence, the impressive performance on specific datasets could be a function of the dataset and not the frameworks themselves. Thus, exploring their efficacies on additional anatomical structures and imaging modalities is also prudent. Initiatives to compile multimodal datasets to train and test frameworks in a challenge-like style such as the Medical Segmentation Decathlon (MSD) could be warranted to ensure that future developments in methodology are sufficiently valid \cite{Antonelli2022}. 
		  
		Nevertheless, longitudinal datasets, be it open-source or in-house, remain scarce. Gathering additional longitudinal data remains the most ideal option, however the aforementioned practical difficulties in data gathering present a significant barrier. In the medium to long term, additional initiatives resembling ADNI could be warranted to gather high-quality, longitudinal, multi-center data for other diseases and disorders benefiting from spatiotemporal shape analyses. Furthermore, these data should be multi-modal, encompassing both imaging and also biomarker data as these have been shown to work synergistically when incorporated into joint frameworks, improving their efficacy. In the meantime, efforts to compile existing data into a large open-source database could be more warranted. This could resemble, for example, the aforementioned MSD. Nonetheless, the impetus to gather and unite such datasets is lacking, especially in the face of general (un)willingness to openly share rare datasets \cite{Tedersoo2021}. 
		

	\subsection{Future outlook and directions}
	
	
		In this review, we focused mainly on the development of LDDMM and DL-based techniques. We did so because these were considered the most versatile for generalizable spatiotemporal shape modeling. This is opposed to alternative methods seeking to model shape changes of specific anatomical structures over time from a mechanistic standpoint. For example, many early works on spatiotemporal shape modeling of tumors attempted to develop models uncovering the underlying mechanistic cause and effects governing their growth \cite{Jarrett2018}. Similar works also exist focusing on cardiac tissue remodeling \cite{Wang2017333333} and bone remodeling \cite{Kameo2020}. Whilst these varied mechanistic models are inherently different, they generally revolve around shape change as a consequence of mechanical and biochemical stimuli or a combination thereof. Thus, these models seek to uncover the underlying formulae governing these interactions and their relationships. This is in contrast with the LDDMM framework, which operates solely from a geometric perspective in uncovering the trajectories of diffeomorphic transformations. In other words, the LDDMM framework does not explicitly consider the underlying physical laws governing the biological processes that lead to the resulting shape changes. Therefore, this approach potentially neglects key information that may affect how reflective the LDDMM approach is in said processes and, therefore, its accuracy. Similarly, DL techniques are opaque, often referred to as black boxes \cite{Castelvecchi2016}. Therein, the model layers, in effect, operate on hidden features uncovered during training processes. These have, in principle, no physical meaning and are not always explainable, engendering issues of trust and validity. A compromise and potential future direction of research is via physics-informed neural networks (PINNs) \cite{Cuomo2022}. Therein, the strengths of DL to process large datasets are utilized to solve underlying physical equations that describe the physics of a system. PINNs are particularly useful even, for example, to uncover underlying dynamics of systems that were previously obscured under high dimensional nonlinear data \cite{Lagergren2020}. Tajdari \textit{et al.} demonstrated the applicability of PINN principles within frameworks to model the longitudinal progression of adolescent idiopathic scoliosis, outperforming traditional methods \cite{Tajdari2021, Tajdari2022}. While their works were mainly concerned with the mechanistic effects of loading on spinal outcomes, their efficacy also lends itself to potential benefits for spatiotemporal shape modeling. Nevertheless, PINNs remain an unexplored avenue and warrant further study. 
		
		Another developing field is utilizing and exploiting causality in the form of causal deep learning. In essence, causality and structural causal models (SCMs) seek to capture and model the chain of causality and inter-variability of multivariate systems \cite{peters2017elements, Pearl2013, pearl2016causal}. This is useful as it enables us to interrogate models to obtain counterfactuals (\textit{i.e.,} if $X$ was different, what is the effect on $Y$? Or more relevantly, 'How would injury $A$ affect bone development of pediatric subject $B$?'). For spatiotemporal shape modeling specifically, this could be used to obtain predictions of shape change over time as a counterfactual from existing data. Traditional methods relied on a system of structural equations with computation graphs, but this precluded the use of higher dimensional data such as images. In recent years, several studies have explored extending SCMs towards being supported by DL (\textit{i.e.,} Deep Structural Causal Models (DSCMs)), enabling the use of hidden features identified via DL \cite{Lore2018, kjsadflkjaslkdjfsdf, heeeeeeeeeeeeeeeeeeeee}. Zhou \textit{et al.} reviewed the synergistic capabilities of generative models and causality, specifically highlighting the applicability of the latter in enhancing the interpretability of generative processes \cite{hehehehehehehhehehehe}. Further works such as by Reinhold \textit{et al.} demonstrated the capability of DSCMs to generate counterfactual brain MRIs of patients with multiple sclerosis \cite{Reinhold2021}. They were able to manipulate demographic and disease covariates and observed their effects on MRI imaging in a novel proof-of-concept. Rasal \textit{et al.} further extended DSCMs towards shape modeling, specifically 3D meshes, demonstrating the extendability and scalability of the principle towards more complex data types \cite{Rasal2022}. Nevertheless, the field is still in its relative infancy, with further developments and refinements in the DSCM framework potentially leading to enhanced efficacy for longitudinal shape modeling. 

        Finally, another potential avenue for development is \textit{via} diffusion models (DMs) \cite{FuestMaGui24, 10081412, Kazerouni2023}. Similar to GANs and AEs described earlier, DMs are a form of generative model where a forward and inverse process are learned to reconstruct input data. In contrast with the aforementioned network architectures, DMs instead function on the principle of noise addition and removal. DMs consist of forward and reverse processes, wherein noise is added onto input data in successive steps and the resulting noise is reversed to reform the input data \cite{DDPM20}. DMs have led to state-of-the-art high resolution visual generative networks \cite{HRIS2021, DMBGIS21}, and existing works have demonstrated the potential for use in a variety of spatiotemporal modeling tasks \cite{ASDMTSSTD2024, NEURIPS2023_8df90a14}. The use of DMs for spatiotemporal shape modeling is still relatively unexplored, with most studies focused on alternative spatiotemporal data (\textit{e.g.,} weather forecasting, traffic patterns, ECG). Nonetheless, Yoon \textit{et. al.} demonstrated the use of a sequence aware diffusion model (SADM) to generate longitudinal medical images \cite{Yoon2023}. Their framework utilized a sequence aware transformer as the conditional module for a diffusion model, demonstrating effective data generation capabilities for longitudinal 3D medical imaging sequences, even with missing data.

		In summary, this paper mainly reviewed the LDDMM framework and DL-based techniques for longitudinal shape modeling. Both achieve their remarkable state-of-the-art performance as they function on similar principles of uncovering the underlying nonlinear spatiotemporal data manifold. Whilst promising, the LDDMM framework is computationally expensive and inefficient due to the exhaustive optimization procedure necessary to calculate smooth and invertible diffeomorphisms. It, nevertheless, demonstrates strong capabilities to establish hierarchical models that differentiate individual and population-level temporal trajectories. Conversely, DL-based techniques are powerful but data-hungry and lack underlying physical meaning. Network architectures have been developed to predict shape changes in anatomical structures. Nevertheless, the underlying data manifolds and spatiotemporal trajectories governing these predictions are obscured by the 'black box' nature of DL architectures. This affects the interpretability of these predictions, especially if the longitudinal trajectories themselves are important. Nevertheless, the capability of DL architectures to identify hidden features from input images and implicitly map the underlying data manifolds denote their importance for spatiotemporal shape modeling. Our review highlights that hybrid techniques that amalgamate both approaches' strengths are more desirable. Furthermore, frameworks incorporating multi-modal data improved generalizability. Thus, further works should not neglect the utility of auxiliary data (\textit{e.g.,} biomarker levels, demographic information, \textit{etc.}). Finally, we theorize that utilizing mechanistic models in a manner similar to PINNs or structured causal frameworks could also further improve the predictive capacities of future spatiotemporal shape models. 
